% Original analytic teaching diagram, 2026-11-12.
% Probability tables on {1,2,3} as the triangle with vertices delta_1=(0,0), delta_3=(60,0),
% delta_2=(30,51.96) (mm). A table (p1,p2,p3) sits at (30 p2 + 60 p3, 51.96 p2).
% Target Pi=(0.6,0.3,0.1) -> (15,15.59).
% Family Q_t=(t,1-2t,t): vertical segment x=30 from (30,0) [t=1/2] to delta_2 [t=0].
% Minimiser of D(Q_t||Pi): t=(4-sqrt6)/5=0.3101, Q=(0.310,0.380,0.310) -> (30,19.73).
% Minimiser of D(Pi||Q_t): t=0.35, Q=(0.35,0.30,0.35) -> (30,15.59), same height as Pi (equal p2).
\definecolor{VB}{HTML}{7188A3}\definecolor{VBF}{HTML}{E5EAF0}
\definecolor{VG}{HTML}{829C8B}\definecolor{VGF}{HTML}{E8EEE9}
\definecolor{VY}{HTML}{B4A16F}\definecolor{VYF}{HTML}{F1ECDD}
\definecolor{VR}{HTML}{A7787D}\definecolor{VRF}{HTML}{F0E5E6}
\definecolor{VP}{HTML}{9588A5}\definecolor{VPF}{HTML}{ECE8F0}
\definecolor{VI}{HTML}{3B4252}\definecolor{VM}{HTML}{6B7280}\definecolor{VL}{HTML}{4C4D4F}
\begin{tikzpicture}[x=1mm,y=1mm,>=Stealth,font=\figurefont\fontsize{8}{10}\selectfont,
lab/.style={text=VM,font=\figurefont\fontsize{7.5}{9}\selectfont}]
\coordinate (D1) at (0,0);
\coordinate (D3) at (60,0);
\coordinate (D2) at (30,51.96);
\draw[VL,line width=.6pt] (D1) -- (D3) -- (D2) -- cycle;
\draw[VL,line width=1.6pt] (30,0) -- (D2);
\draw[VB,line width=.6pt,dashed] (15,15.59) -- (30,15.59);
\fill[VB] (15,15.59) circle (1.2);
\fill[VY] (30,19.73) circle (1.2);
\fill[VP] (30,15.59) circle (1.2);
\foreach \p in {(D1),(D2),(D3)} {\fill[VI] \p circle (0.7);}
\node[text=VI,anchor=north east] at (D1) {$\delta_1$};
\node[text=VI,anchor=north west] at (D3) {$\delta_3$};
\node[text=VI,anchor=south] at (D2) {$\delta_2$};
\node[text=VB,anchor=east] at (14,15.59) {$\Pi$};
\node[text=VY,anchor=west] at (31.5,20.6) {$t=0.310$};
\node[text=VP,anchor=west] at (31.5,14.4) {$t=0.35$};
\node[lab,anchor=north] at (30,-0.8) {$t=\tfrac12$};
\node[lab,anchor=west] at (31.5,38) {$Q_t$};
\end{tikzpicture}
